\begin{recipe}
[% 
    preparationtime = {\unit[20]{min} + \unit[1]{h} Ruhezeit},
    portion = {\portion{3}},
    %source = {},
    calory = {292kcal | 10g Protein | 60g KH | 1g Fett},
    price = {0,27€},
    created = {11.03.2026},
    rating = {3}
]
{Pasta Orecchiette}

\graph{% pictures
    big=pic/grundlagen/pasta-orecchiette
}

\introduction{%
Ein einfacher Pastateig nur aus Hartweizengrieß und Wasser. Die geformten Orecchiette („Öhrchen“) halten Sauce besonders gut und sind überraschend leicht selbst zu machen.
}

\ingredients{
    250 g & Hartweizengrieß\index{Getreide, Pasta \& Brot!Hartweizengrieß}\\
    125 ml & Wasser\index{Getränke!Wasser}\\
    \nicefrac{1}{2} TL & Kurkuma\index{Kräuter \& Gewürze!Kurkuma}\\
    & Salz\index{Kräuter \& Gewürze!Salz}
}

\preparation{%
    \step%
    Salz im Wasser auflösen. Hartweizengrieß und Kurkuma in eine Schüssel geben und das Wasser hinzufügen.
    
    \step%
    Alles zu einem festen Teig verkneten (ca. 8--10 Minuten), bis er glatt und elastisch ist. Beim eindrücken mit dem Finger, sollte der Teig zurückspringen.
    
    \step%
    Teig abdecken und etwa 1 Stunde ruhen lassen.\\
    
    \step%
    Den Teig zu daumenbreiten Rollen formen und in kleine Stücke schneiden.
    
    \step%
    Mit einem Messer über die Stücke ziehen, sodass sie sich leicht aufrollen.
    
    \step%
    Die Stücke mit dem Finger nach außen stülpen, sodass die typische Orecchiette-Form entsteht.
    
    \step%
    Pasta in gut gesalzenem Wasser etwa 4--6 Minuten kochen, bis sie nach oben schwimmen und al dente ist.
}

\suggestion{
    Das Salz sollte im Wasser volständig aufgelöst sein, da sonst die Kristalle beim ausrollen die Strukur kaputt machen kann.
}

\end{recipe}